\section{Una perspectiva unificada de Flow Matching y los modelos de difusión}

\subsection{Establecimiento de la equivalencia}

En esta parte final de la clase, el docente reveló las conexiones profundas entre Flow Matching y los modelos de difusión. Cuando seleccionamos $\mu_t$ y $\sigma_t$ correspondientes a los esquemas VP/VE:

\begin{enumerate}
    \item La trayectoria de probabilidad condicional de Flow Matching $p_t(x | x_1) = \mathcal{N}(\mu_t x_1, \sigma_t^2 I)$ es idéntica a la distribución del salto hacia adelante del modelo de difusión (solo con dirección temporal opuesta).
    \item Existe una correspondencia exacta entre el campo vectorial condicional de Flow Matching y el término de deriva en la ODE de Probability Flow.
    \item Ambos objetivos de entrenamiento son equivalentes bajo una parametrización adecuada y ponderación de pérdidas apropiada.
\end{enumerate}

\begin{importantbox}{Perspectiva unificada: dos caminos hacia la misma montaña}
\textbf{Trayectoria del modelo de difusión}: Definir el proceso hacia adelante (inyección de ruido) $\to$ derivar la SDE inversa $\to$ obtener la ODE de Probability Flow $\to$ entrenar la función score.

\textbf{Trayectoria de Flow Matching}: Definir el mapeo de flujo condicional $\to$ derivar el campo vectorial condicional $\to$ entrenar una red que prediga el campo vectorial.

Ambos caminos llegan finalmente al mismo objetivo: una ODE capaz de mapear ruido a datos. La diferencia radica en:
\begin{itemize}
    \item El modelo de difusión deriva el campo vectorial partiendo de la trayectoria de densidad de probabilidad.
    \item Flow Matching deriva la trayectoria de densidad de probabilidad partiendo del mapeo de flujo.
    \item El orden de derivación es inverso, pero los resultados son equivalentes.
\end{itemize}
\end{importantbox}

\subsection{Ventajas de Flow Matching}

Aunque matemáticamente equivalentes, el marco de Flow Matching ofrece las siguientes ventajas prácticas:

\begin{table}[H]
\centering
\begin{tabular}{lcc}
\toprule
\textbf{Característica} & \textbf{Modelo de difusión} & \textbf{Flow Matching} \\
\midrule
Distribución de referencia & Debe ser gaussiana estándar & Puede ser cualquier distribución \\
Objetivo de entrenamiento & Predicción de ruido / Predicción de score & Predicción de velocidad \\
Derivación teórica & Requiere SDE/Fokker-Planck & Define directamente la ODE \\
Esquema de ruido & Requiere diseño cuidadoso & Selección libre de $\mu_t$, $\sigma_t$ \\
Trayectoria recta & No natural (trayectorias curvas) & Natural (trayectoria OT) \\
\bottomrule
\end{tabular}
\caption{Comparación de características entre modelos de difusión y Flow Matching}
\end{table}

\begin{knowledgebox}{Por qué los modelos generativos modernos prefieren Flow Matching}
Modelos de generación de imágenes más recientes como Stable Diffusion 3 (Stability AI) y Flux (Black Forest Labs) han adoptado el marco de Flow Matching / Rectified Flow. Las razones principales incluyen:
\begin{enumerate}
    \item \textbf{Trayectorias más rectas}: Las trayectorias generadas por la trayectoria OT se acercan más a líneas rectas, lo que mejora el rendimiento con pocas muestras.
    \item \textbf{Entrenamiento más sencillo}: El objetivo de predicción de velocidad es más intuitivo que la predicción de ruido o score.
    \item \textbf{Diseño más flexible}: Permite seleccionar libremente la distribución de referencia y la parametrización de la trayectoria.
    \item \textbf{Fundamento teórico superior}: Conexión natural con la teoría de transporte óptimo (optimal transport).
\end{enumerate}
\end{knowledgebox}

\subsection{Resumen de la sección}

Flow Matching y los modelos de difusión son matemáticamente equivalentes, pero Flow Matching proporciona un marco teórico más conciso y un espacio de diseño de modelos más flexible. En la práctica, la estructura de trayectoria óptima derivada de trayectorias rectas (trayectoria OT) hace que Flow Matching sea el paradigma de entrenamiento preferido para los modelos generativos actuales.
